% ============================================================
% 1. INTRODUCCIÓN
% Base: secciones A, C, D y F del anteproyecto. Las directrices exigen
% los objetivos APROBADOS: van literales, no se reescriben.
% TODO: pasar a pasado lo que el anteproyecto dejaba en futuro.
% ============================================================
\chapter{Introducción}
\label{cap:introduccion}

% ------------------------------------------------------------
\section{Planteamiento del problema}
\label{sec:planteamiento}

En holografía digital el sensor no registra la imagen del objeto sino un patrón
de interferencia: para recuperarla hay que modelar numéricamente cómo viajó la
luz hasta el plano del objeto~\cite{buitrago2022}. Dentro de la teoría escalar
de la difracción ese cálculo ya tiene solución cerrada: la integral de
Rayleigh--Sommerfeld y el espectro angular (ASM) son la misma solución escrita
en el dominio espacial y en el de las frecuencias espaciales, y ninguna aproxima
el ángulo de propagación \cite{goodman2005, ersoy2006}. Lo que las separa es el
costo: el ASM se calcula con la transformada rápida de Fourier (FFT) en
$\mathcal{O}(N^2\log N)$ operaciones frente a las $\mathcal{O}(N^4)$ de la suma
directa, y por eso es el método más usado~\cite{buitrago2020}. La contrapartida
es que la exactitud deja de depender de la formulación y pasa a depender del
muestreo con que se discretiza.

La FFT impone sus propias reglas de muestreo: el espaciado de los puntos en
frecuencia queda fijado por el tamaño de la ventana de entrada, y el plano de
salida hereda el tamaño del de partida \cite{goodman2005}. Cuando esas reglas no
alcanzan a describir la propagación aparece \emph{aliasing}: frecuencias altas
que se confunden con bajas y contaminan la imagen. Las soluciones conocidas
pagan un precio: el método de banda limitada (BLAS) descarta las frecuencias
problemáticas y pierde detalle a distancias largas \cite{matsushima2009};
ampliar la ventana de entrada corrige el muestreo pero dispara memoria y tiempo
de cómputo \cite{yu2012}. Una alternativa reciente reemplaza la FFT por un
producto de matrices ---el método de espectro angular por producto matricial
(MPASM)~\cite{zhao2020}---: el muestreo deja de estar impuesto por el algoritmo
y se elige mediante dos coeficientes de control. $K_f$ comprime el intervalo de
muestreo en frecuencia, concentrando los puntos donde la función de
transferencia sí cumple la condición de muestreo; $s$ multiplica el número de
puntos en frecuencia y amplía el ancho de banda efectivo a distancias largas, a
costa de más operaciones.

La microscopía holográfica digital sin lentes (DLHM) es un escenario donde ese
control resulta pertinente. Al eliminar el objetivo del microscopio, la muestra
se ilumina con una onda esférica proveniente de un orificio muy pequeño
(\emph{pinhole}) y el sensor registra directamente el patrón de difracción, de
modo que la resolución queda determinada por la geometría del
montaje~\cite{lopera2024}. Dos rasgos de esa geometría tensionan al ASM
convencional: el paso del sensor deja el contenido frecuencial del campo cerca
del límite de Nyquist, y la iluminación divergente introduce una magnificación
que sitúa el patrón registrado en una escala distinta a la del objeto, algo que
la FFT no puede representar. El MPASM ofrece justamente ese control, pero solo
había sido evaluado en simulaciones con haces gaussianos e iluminación de onda
plana: no se había reportado su aplicación a hologramas experimentales ni
existían criterios publicados para elegir sus parámetros en condiciones reales
de registro.

\vspace{0.2cm}
\noindent\textbf{Pregunta de investigación:} ¿Qué características de un
propagador con control explícito del muestreo, como el MPASM, permiten mejorar
la calidad de las reconstrucciones en microscopía holográfica digital sin lentes
respecto de los propagadores empleados habitualmente en la técnica, y bajo qué
condiciones de registro resulta ventajoso?

% ------------------------------------------------------------
\section{Objetivos}
\label{sec:objetivos}

% Literales del anteproyecto aprobado. No editar.
\subsection*{Objetivo general}

Aplicar propagadores de campo óptico basados en el producto matricial, con
control flexible del muestreo espacial y frecuencial, a la reconstrucción de
hologramas en microscopía holográfica digital sin lentes.

\subsection*{Objetivos específicos}

\begin{enumerate}[label=\arabic*.]
  \item Implementar en un marco computacional unificado los métodos de
  propagación FFT-ASM, BLAS y MPASM, verificando su correcta ejecución mediante
  casos analíticos de referencia.

  \item Validar el desempeño de los métodos implementados bajo las condiciones
  de muestreo propias de DLHM, determinando los rangos de distancia, apertura
  numérica y coeficientes de control ($s$ y $K_f$, definidos en la
  Sección~\ref{sec:planteamiento}) en los que cada método mantiene una
  precisión aceptable, y contrastando la formulación matricial con la de la
  transformada Z \emph{chirp} (CZT), que permite un control de muestreo
  análogo.

  \item Comparar las reconstrucciones obtenidas con los propagadores
  implementados frente al método de espectro angular con iluminación de onda
  esférica, empleado actualmente en el montaje DLHM del laboratorio, mediante
  métricas cuantitativas de calidad sobre hologramas experimentales de muestras
  de resolución conocida.
\end{enumerate}

% ------------------------------------------------------------
\section{Justificación}

El MPASM aporta un grado de libertad que los métodos convencionales no tienen:
elegir de forma independiente cómo se muestrea el campo en el espacio y en la
frecuencia. Ese control encaja con lo que exige DLHM, donde la magnificación
geométrica y las condiciones de registro imponen restricciones que el ASM
basado en FFT no puede acomodar. Verificar si esa flexibilidad se traduce en
mejoras medibles sobre la reconstrucción de hologramas experimentales es una
pregunta que la literatura aún no respondía, ya que el método solo había sido
evaluado en simulaciones; y contar con criterios explícitos para elegir los
parámetros de muestreo evita el ajuste por ensayo y error que exige usar estos
propagadores. A esto se suma que su formulación matricial se parece
estructuralmente a la transformada Z \emph{chirp} (CZT), un algoritmo que
también permite calcular espectros con muestreo arbitrario; observar esa
relación durante la caracterización numérica ayuda a situar el método entre los
propagadores existentes.

% ------------------------------------------------------------
\section{Alcance}

El trabajo se limita al régimen de difracción escalar y a hologramas en línea
registrados en un montaje DLHM existente; no incluye instrumentación nueva, el
régimen vectorial ni algoritmos de recuperación de fase.

\pendiente{productos entregados: el paquete de software (nombre, licencia,
repositorio), los hologramas con sus reconstrucciones y los mapas de validez.
Decirlo en pasado y con lo que de verdad se entregó.}

% ------------------------------------------------------------
\section{Estructura del documento}

\pendiente{un párrafo que recorra los capítulos: el
Capítulo~\ref{cap:marco} presenta el marco teórico y los antecedentes; el
Capítulo~\ref{cap:metodos}, la implementación, los barridos y el montaje; el
Capítulo~\ref{cap:resultados}, los resultados por objetivo; y el
Capítulo~\ref{cap:conclusiones}, las conclusiones y el trabajo futuro.}
